\sectionkicker{Theme synthesis}
\subsection*{横断比較 — Live environmentを何で置き換えるか}
\addcontentsline{toc}{subsection}{横断比較 — Live environmentを何で置き換えるか}
DynaWebとSYNTHAGENTはどちらもreal environment依存を下げるが、代替しているものが異なる。world modelによるrolloutと、task・tool・user・rewardの合成を分けると設計上のtrade-offが見える。
\medskip
\begin{center}
\small
\renewcommand{\arraystretch}{1.16}
\begin{tabularx}{\linewidth}{@{}p{0.20\linewidth}X@{}}
\toprule
観察軸 & 7月の一次資料・論文から読めること \\
\midrule
環境の表現 & DynaWebはweb page representationの遷移を予測するlearned world modelを使い、その中でmodel-based RL rolloutを行う。live webそのものではなく学習した環境modelが訓練surfaceになる。 \cite{src-85bc3e4f7b39} \\
task / toolの合成 & SYNTHAGENTはsynthetic tasks、mock tools、LLM user simulator、rubric rewardsを組み合わせ、taskとtool interactionの側を生成する。world transitionの学習とは異なる合成単位である。 \cite{src-01d8ab4a2c6f} \\
運用上の狙い & 両手法はreal environmentのcost、privacy、instabilityを下げ得る方向を示す一方、benchmark結果や改善幅は著者評価として読む必要がある。production環境での一般的優位性へ直接変換しない。 \cite{src-85bc3e4f7b39,src-01d8ab4a2c6f} \\
simulation gap & 学習済みworld modelもsynthetic tool environmentも、dynamic・adversarialなlive systemのfailureを完全には再現しない可能性がある。さらにtask locatorは後日改訂版を含むため、1月のchronologyはoriginal submissionへ固定する。 \cite{src-85bc3e4f7b39,src-01d8ab4a2c6f} \\
\bottomrule
\end{tabularx}
\renewcommand{\arraystretch}{1.0}
\end{center}
\normalsize
\begin{claimboundary}[この比較の境界]
この表は本号で既に検証・参照している一次資料と論文を横断比較の形へ再配置したものである。vendor / project / author が報告した性能値や優位性は独立再現済みの一般則へ変換せず、本文とSource Notesの帰属境界をそのまま維持する。
\end{claimboundary}
